\chapter{Conclusion and Future Work}
\section{Answers to the primary question, RQ1 and RQ2}

The primary question is answered within the controlled scope. The implemented workflow can preserve
identity, reporting-period, missing-state, provenance, verification and approval records from intake
to report export. Engineering checks and repository artefacts establish that these boundaries exist in
the evaluated local version. They do not establish operational usefulness or production readiness
\citep{hevner2004design,peffers2007dsrm}.

For RQ1, the strongest result is traceable execution. D0 decisions can be followed from a fictional
candidate and its evidence to the verifier state and report consequence, and the complete outcome list
was identical across three repeats. The observed-zero, contradiction, stale-period and untrusted
examples show what the evidence rules mean in practice. This answers how the implemented controls
keep a report's meaning intact within D0 \citep{gao2023alce,gao2023rarr,pineau2021reproducibility}.

For RQ2, separate verification changed claim admission under equal inputs. C1 emitted six candidates
that D0 labelled unsupported, while C2 held all six and retained the five supported claims. Precision
changed from 0.455 to 1.000, the unsupported-claim rate from 0.545 to 0.000, and recall remained 1.000.
This is an observed verification-gate effect on fourteen designed cases. It is not evidence that
multiple agents are generally superior or that the same rates will hold for real companies
\citep{guo2024multiagent,kapoor2023leakage}.

The dissertation's conclusion is therefore narrower and stronger than the original research
contract. It demonstrates an evidence-first implementation and one controlled verification effect.
Manual comparison, participant review, live source use and business benefit are outside the active
questions because the required observations were unavailable. Their absence does not make the
executed result incomplete; it sets the boundary for future evaluation
\citep{hevner2004design,pineau2021reproducibility}.

\section{Demonstrated contribution}

The main contribution is a reporting workflow that makes source admission, claim support, conflicts,
missing information and release authority explicit. Separating the roles gives production,
verification and approval different permissions and saved hand-offs. The design is inspectable
because the saved records link each decision to a source, rule, version and output
\citep{buneman2001provenance,pineau2021reproducibility}.

The empirical contribution is the paired D0 comparison. It reports aggregate metrics and exact case
examples, including a correct accepted zero and three different reasons for holding a claim. The result
supports the verification contract on the fixture rather than a new algorithm, a persona effect or a
claim of production performance \citep{gao2023alce,gao2023rarr}.

The practical contribution is an executable pilot design. It names the existing XLSX input and
Next.js review interface, assigns staff responsibilities, defines cost and time records, specifies
quality and operational measures, and sets go/no-go and stop rules. Time saving and cost reduction
remain hypotheses until that pilot is authorised and run
\citep{peffers2007dsrm,nist2023airmf}.

\section{Sequenced future evidence}

The first next step is the prospective business pilot in Section~6.7. It should freeze the data,
source window, reference decisions, completion rule and acceptance criteria before collecting paired
manual and workflow observations. The pilot must record active and elapsed time, reviewer work,
supported-claim precision, evidence coverage, false acceptance, false rejection, source failures,
completion and cost per report. A pilot result should not be used to tune a sealed final evaluation
set \citep{kapoor2023leakage,pineau2021reproducibility}.

A separate live-source study can then compare approved official connectors with bounded assisted
discovery using identical companies, dates, source limits and output requirements. Any external model
route requires current provider controls, approved data classes and recorded usage. The outstanding
testing listed at the end of Section~6.7 must be finished before remote or multi-user operation is
considered \citep{nist2023airmf,cddo2023genai}.

A managed platform such as Grok Bot should then be tested against the custom orchestration layer on
the terms set out in Section~6.4. The question is whether a bought platform can reproduce the same
evidence-admission, verification and approval contract for less setup and maintenance effort, measured
under equivalent inputs rather than assumed from its feature list
\citep{spacexai2026grokbot,guo2024multiagent,pineau2021reproducibility}.

The full evidence roadmap is Table~\ref{tab:conc-evidence-roadmap} in
Appendix~\ref{app:supplementary-exhibits}. Each future stage
requires an explicit authority and stop decision. None is presented as a result of this dissertation
\citep{nist2023airmf,pineau2021reproducibility}.
